\noindent Throughout: $q = 1-p$, $n = $ number of trials, $k = $ number of successes,
and $\ff{x}{r} = x(x-1)\cdots(x-r+1)$ is the $r$th falling factorial.

\section{The Binomial Distribution}

Let $S_n$ count the number of successes in Bernoulli trials $Z_i$, $i = 1,\dots,n$, with $Z_i \iid Z \sim \bern(p)$.

\begin{definition}[Binomial distribution]
The random variable $X$ has a \emph{binomial distribution on $n$ trials with probability of success per trial $p$}, denoted
\[
X \sim \bin(n,p), \qquad n \in \N,\ 0 \le p \le 1,
\]
if and only if
\[
X \eqd S_n = Z_1 + \cdots + Z_n \quad \text{with } Z_1,\dots,Z_n \iid Z \sim \bern(p).
\]
\end{definition}

For $X \sim \bin(n,p)$, by linearity of expectation:
\[
\E X = \E S_n = \E Z_1 + \cdots + \E Z_n = n\,\E Z = np.
\]
\fix{The handwritten notes say ``for $X \sim \bern(p)$''. It should be $X \sim \bin(n,p)$, because a single Bernoulli has mean $p$, not $np$.}

\begin{proposition}[Sums of independent binomials]
If $X \sim \bin(m,p)$, $Y \sim \bin(n,p)$ and $X \indep Y$, then
\[
X + Y \sim \bin(m+n,\,p).
\]
\end{proposition}

\begin{corollary}
If $X_i \overset{\text{indep}}{\sim} \bin(n_i, p)$ for $i = 1,\dots,k$, then
\[
\sum_{i=1}^{k} X_i \sim \bin\!\Big(\sum_{i=1}^{k} n_i,\ p\Big).
\]
\end{corollary}
\begin{remark}
The success probability $p$ has to be the \emph{same} for every summand.
\end{remark}

\subsection{Deriving the pmf}

Any single outcome of the $n$-vector of $Z$-trials is one of the $2^n$ sequences of $0$'s and $1$'s in $\R^n$.
So the effective sample space for $\mathbf{Z} = (Z_1,\dots,Z_n)$ is the set of $2^n$ corners
\[
\{0,1\}^n, \qquad \#\{0,1\}^n = 2^n,
\]
of the unit cube $[0,1]^n \subset \R^n$.

\begin{example}
For $n = 2$: $\{0,1\}^2 = \{(0,0),(0,1),(1,0),(1,1)\}$.
\end{example}

Since the trials are independent,
\[
P(\mathbf{Z} = \mathbf{z}) = \prod_{i=1}^{n} P(Z_i = z_i) \qquad \forall\, \mathbf{z} \in \{0,1\}^n.
\]

\begin{example}
$P\big(\mathbf{Z} = (1,1,0)\big) = P(Z_1 = 1)\,P(Z_2 = 1)\,P(Z_3 = 0) = p \cdot p \cdot q = p^2 q.$
\end{example}

Individually, $P(Z_i = z_i) = p^{z_i} q^{1-z_i}$ for $z_i \in \{0,1\}$, $i = 1,\dots,n$. Therefore
\[
P(\mathbf{Z} = \mathbf{z}) = p^{\sum z_i}\, q^{\,n - \sum z_i} \qquad \forall\, \mathbf{z} \in \{0,1\}^n.
\]

For each $k = 0,1,\dots,n$ define the set of corners with exactly $k$ ones:
\[
C_k^n = \Big\{ \mathbf{z} \in \{0,1\}^n \;\Big|\; \textstyle\sum z_i = k \Big\}.
\]
\fix{The notes have $k = 1,\dots,n$. It should start at $k = 0$, because zero successes is a possible outcome.}

\begin{example}
For $n = 3$, $k = 2$: $C_2^3 = \{011,\,101,\,110\}$.
\end{example}

Then $S_n = k \iff \mathbf{Z} \in C_k^n$, so
\[
P(S_n = k) = \sum_{\mathbf{z} \in C_k^n} P(\mathbf{Z} = \mathbf{z})
= \sum_{\mathbf{z} \in C_k^n} p^{\sum z_i} q^{\,n-\sum z_i}
= \#C_k^n \cdot p^k q^{n-k}.
\]
Every term in the sum is the same, $p^k q^{n-k}$, which is why the sum collapses to a count.

\begin{example}
$P(S_3 = 2) = P(011) + P(101) + P(110) = 3p^2 q$.
\end{example}
\fix{The notes list $P(001) + P(110) + P(101)$. But $001$ has only one success, so it should be $011$.}

Writing $\#C_k^n = \binom{n}{k}$ (``$n$ choose $k$''):

\begin{definition}[Binomial pmf]
For $X \sim \bin(n,p)$,
\[
P(X = k) = \binom{n}{k} p^k q^{n-k}, \qquad k = 0,1,\dots,n.
\]
\end{definition}

\begin{corollary}
For any $n \in \N$ and $0 \le p \le 1$,
\[
\sum_{k=0}^{n} \binom{n}{k} p^k q^{n-k} = \sum_{k=0}^{n} P(S_n = k) = 1.
\]
\end{corollary}

\begin{corollary}
\[
\binom{n}{k} = \frac{\ff{n}{k}}{k!} = \frac{n(n-1)\cdots(n-k+1)}{k!}.
\]
\end{corollary}

\begin{proposition}[Behaviour in $n$ for fixed $k$]
For fixed $k$:
\begin{enumerate}[label=(\roman*)]
  \item for $n > k/p$, $P(S_n = k)$ is strictly decreasing in $n$;
  \item $\displaystyle\lim_{n\to\infty} P(S_n = k) = 0$.
\end{enumerate}
\end{proposition}

\subsection{Factorial moments}

Computing $\E X^2, \E X^3, \dots$ directly from the pmf is messy. Falling factorials make it easy:
\[
\ff{X}{r} = X(X-1)\cdots(X-r+1).
\]

\begin{proposition}
For $X \sim \bin(n,p)$ and $r = 1,\dots,n$,
\[
\E \ff{X}{r} = \sum_{k=0}^{n} \ff{k}{r} \binom{n}{k} p^k q^{n-k} = \ff{n}{r}\, p^r,
\]
and $\E \ff{X}{r} = 0$ for $r > n$.
\end{proposition}

\begin{example}[Second moment and variance]
$\ff{X}{2} = X(X-1) = X^2 - X$, so
\[
\ff{n}{2} p^2 = \E \ff{X}{2} = \E X^2 - \E X = \E X^2 - np.
\]
Hence
\[
\E X^2 = n(n-1)p^2 + np, \qquad
\operatorname{var} X = \E X^2 - (\E X)^2 = n(n-1)p^2 + np - n^2p^2 = npq.
\]
\end{example}
\fix{The notes conclude $\E \ff{X}{2} = n^2 p^2$. That's wrong: $\ff{n}{2} = n(n-1)$, not $n^2$. If you use $n^2$, you get $\E X^2 = n^2p^2 + np$ and $\operatorname{var} X = np$, both of which are wrong.}

\section{The Negative-Binomial Distribution}

Let $T_k$ count the number of trials required until (and including) the $k$th success:
\[
T_k = \min\{ n \in \N \mid S_n \ge k \}.
\]

\begin{center}
\color{Gray}\itshape [Notes incomplete --- section to be continued.]
\end{center}
